\section{What remains uncertain}
\label{sec:unc}

Everything below is a limitation of the deliverable as it stands, with the number that
measures it. None of it is a reason not to use the method; all of it is what a reader should
know before relying on a particular number.

\subsection{Two things that are not proven}

\paragraph{$\mathrm{est}(\mathbf{R}) \ge \max_{ab}|T_{ab}|$ is not a theorem.} The selection
rule compares the proved bound against $\mathrm{tol}\cdot\mathrm{est}(\mathbf{R})$, and the
statement ``relative error $\mathrm{tol}$'' needs $\mathrm{est} \le \max_{ab}|T_{ab}|$,
which is false: measured over the $427$ references, $\mathrm{est}/\max|G_{\mathrm{ref}}|$
runs over $[1.10, 6.34]$ on the cubes and reaches $135$ on the slender needle. It
over-states, which is the safe direction --- the rule certifies
$\mathrm{tol}\cdot\mathrm{est}/\max|T|$, i.e.\ $1.1$ to $6.3$ times $\mathrm{tol}$ --- but
it is a measurement and not a proof, and if a shape existed for which $\mathrm{est}$
under-stated the true magnitude the selected $L$ would be too small. The cross-scale round
found such a case: on the $16{:}1{:}1$ rod at one fine cell of gap
$\mathrm{est}/\max|G|$ is $0.50$ at $f = 1$ and $0.25$ at $f = 0.37$
(Section~\ref{sec:xs:open}). The whole box does not converge at that offset, so no $L$ was
selected from the under-stated scale, but the hypothesis is now known to fail on a shape
Gila produces. The a posteriori
certificate \eqref{eq:cert:post} discharges the hypothesis entirely at no cost and evaluates
to $4.5\times10^{-14}$ over the same references; the proven lower scale
$\mathrm{est}_{\mathrm{lo}}$ exists but is vacuous at $\lambda/4$
(Section~\ref{sec:bnd:cert:post}) and fails at the nearest unequal offsets for ratios
$\ge4$ on two or three axes (Remark~\ref{rem:xs:est}).

\paragraph{The rounding error is measured, not bounded.} Section~\ref{sec:bnd:cert} states
the number: $\le 16.4\,\varepsilon\max_{ab}|T_{ab}|$ over sixty random cases spanning aspect
ratios $1$ to $1020$, with an $l$-sum amplification of at most $3.14\times10^{3}$. Nothing
in this document bounds it. Related, and also measurements rather than theorems: the upward
recurrence for $h_l$ is accurate to $3\times10^{-15}$ for $l \le 40$,
$|z|\in[0.3,300]$, $\operatorname{Im}z\in[0,30]$, but that is $72$ measured rows and not a
proof --- which is why the shipped bound uses the \emph{certified} majorant
$\mathrm{hb}(l,z)$ of Theorem~\ref{thm:bnd:h} rather than $|h_l|$ itself, at a cost of
$1.001$--$1.008$ in per-offset work on cubes and $1.2$ on the slender cell. Miller's
downward recurrence returns $j_l$ to $6.8\times10^{-15}$; the odd-$l$ radial moments use
$W_l \le \sqrt{W_{l-1}W_{l+1}}$, exact in the even case and a genuine bound in the odd; and
the monotonicity in $|\mathbf{R}|$ that the threshold table assumes was checked against a
direct evaluation at $4000$ radii per (shape, frequency), which is a check and not a proof.

\subsection{Where the accuracy statement has to be qualified}

\paragraph{The small part of a complex entry.} For real $f$ the tensor splits into two
independent real sums, $\operatorname{Re}T$ from $y_l$ and $\operatorname{Im}T$ from $j_l$,
and each part keeps its own relative accuracy. For complex $f$ it is one complex sum, and
then a part carrying less than about $10^{-3}$ of the entry's own modulus is accurate only
to $\varepsilon\max_{ab}|T_{ab}|$ in absolute terms, i.e.\ to a \emph{relative} error of
about $10^{-16}$ divided by that fraction. This is visible in the Re/Im columns of every
verification table at $10^{-13}$ to $9\times10^{-12}$, and it is a property of the
formulation, not of the expansion: the independent Cartesian route has it too, and none of
the seven families in the registry avoids it. It is not fixable in \code{Float64} without
splitting the complex-frequency sum by magnitude. \textbf{Entries themselves --- and, at
real $f$, parts of entries --- meet $10^{-14}$; parts of entries at complex $f$ that carry
less than $1\,\%$ of their entry do not.}

\paragraph{Aspect ratios were verified only where references exist.} The library is generic
in three independent rational edge lengths and so is the bound, but a reference is an
expensive object. What exists: the four production shapes at five frequencies ($427$
tensors), the twelve audit shapes at aspect ratios $1$ to $1020$ at four frequencies
($288$ tensors, of which four rows are limited by the reference's own $3.6\times10^{-15}$
floor at an aspect-$678$ needle two cells apart), and sixty random cases. \emph{No claim is
made about a shape outside that set beyond what the shape-generic bound gives.}

\paragraph{Cells above $\lambda/4$ are refused or expensive, not wrong.} Above
$|k|r_{\mathrm{d}} \approx 3.2$ the $k$-series inside $j_l$ needs more terms than a
$\lambda/4$ cell does; the library now computes how many and builds the table to match, and
above $N_{\mathrm{cap}} = 40$ ($|k|r_{\mathrm{d}} \approx 22$) it raises an error rather than
truncating. But $N_{\mathrm{cap}}$ is a policy, not a theorem: nothing says the expansion is
the right tool there, and a $\lambda/2$ cube needs $N = 16$ and $171$~s of table, a
$\lambda$ cube $N = 24$ and more. There is also a reversal worth knowing: on a $\lambda/2$
cell the two-cell offset does not converge on either expansion route and falls through to
the $k$-series, which answers it to $5.6\times10^{-17}$ --- \emph{on a coarse cell the near
field is the accurate part and the far field is the expensive one}, the opposite of the
usual expectation.

\subsection{Cost and engineering}

\paragraph{The memory of the table build is measured on four shapes.} The in-place
rewrite of Section~\ref{sec:use:cost} brought the cold build from $21$--$25$~GB to
$0.66$--$0.74$~GB on the four production shapes and was checked bit-identical on $24$
tables. The twelve audit shapes of aspect ratio up to $1020$ were not rebuilt after the
rewrite; their tables are the same code path at the same $L$ and $n_{\max}$, so nothing
suggests a different outcome, but it is inferred and not measured.

\paragraph{The cost discontinuity at the $k$-series boundary.} One offset of the slender
needle costs $1.9$ to $9.5$~s and its lattice neighbour one lattice step further out costs
$0.13$~ms: a factor of $10^{5}$ across one step. It is a property of the routing rule, not of
the machine, and the route has no cost cap --- it raises only when $200$ series terms are
not enough. A $\lambda/8$ plate at $N = 28$ still takes tens of seconds per offset. It is
paid once per (shape, frequency) and cached, and it is $67$ offsets of half a million, but a
user meeting it cold on a new shape will wait two minutes on eight threads.

\paragraph{Precision above $192$ bits buys nothing.} The geometry table is stored at
\code{TABPRC = 192}, so a \code{BigFloat(256)} result advertises $77$ digits and carries
$58$. The type is generic; the accuracy is capped.

\paragraph{No GPU path.} \code{farBlock!} is CPU-threaded and bitwise deterministic in the
thread count. Nothing here was written for or tested on a GPU, and the machine used had no
CUDA device.

\paragraph{Wiring into \code{GlaVacOprMem} is a later task.} \code{farBlock!} fills a
Gila-layout \code{egoToe} and \code{bench.jl} drives the real build stage by stage through
\code{GilaElectromagnetics.GilaVacuum}, verifying that the staged ``before'' path is
bit-identical to \code{GlaVacOprMem}'s own output and that the ``after'' path differs from it
only by Gila's own quadrature error. But nothing under \code{src/} was touched, by
instruction. Replacing the \code{egoFunInn!} loop in \code{genEgoCrcSlf!} is the next task,
and it needs a decision on where the shape tables live and who builds them.

\subsection{Unequal cells}

Section~\ref{sec:xs:open} lists what the cross-scale extension leaves open; the items that
qualify a number in this document are these.

\paragraph{The certificates of unequal pairs are loose where $\mathrm{est}$ is far off,
and $0.13\,\%$ of a realistic block is certified only to $1$--$4\,\mathrm{tol}$.} Across
the thin axis of the slender pair $(1/32,1/32,1/512)$ against $(1/8,1/8,1/512)$ the
tensor is $10^{2}$--$10^{4}$ below the pointwise scale, so the route-3 certificate is
$2600$--$8900$ times the actual error there ($1.5\times10^{-14}$ to $8.1\times10^{-13}$ of
$\max|G|$ against errors of $10^{-17}$). On the $\lambda/2$ block $1800$ offsets at
$\rho_{\mathrm{whl}} = 0.55$--$0.60$ meet the bound at $L = 56$ only to $1$--$4$ times
$\mathrm{tol}\cdot\max|G|$; where a reference exists they are correct to
$3\times10^{-15}$, but that is what is proven.

\paragraph{The verification of unequal cells covers ratios $2$, $4$, $8$, $16$ on one to
three axes against a $\lambda/32$ fine cell, two slender pairs, and three frequencies,
and the adversarial round added ratios $3$ and $6$, a pair coarse on a different axis in
each cell, coarse cells of $\lambda/2$ and $\lambda$, two complex frequencies and
mixed-sign offsets (Section~\ref{sec:xs:adv}).} What that round did not turn into a
verified case: non-integer rational ratios, which the library refuses rather than computes
--- a three-line change makes all three routes work on them and was measured but not
shipped; coarse cells above $\lambda$ and ratios above $32$; and \code{Float32} anywhere in
the near band, which returns \code{NaN} rather than a poor answer (D14). Route 3 on unequal
panels was verified through $103$ off-lattice reference records and never through a block,
since no block offset reaches it. The routing itself is no longer a sample: it is exhaustive
over the $125\,256$ separated offsets of the four test blocks, and \code{farBlockX!} is
bitwise deterministic in the thread count.

\paragraph{Two costs are engineering, not mathematics.} The exact-rational offset vector
costs as much to classify as the far field costs to compute ($4.2$~s against $4.4$~s
single-threaded on the $\lambda/2$ block, and about $1$~GB of \code{BigInt}); an
integer-lattice interface would remove it. The table writer still appends without a
cross-process lock (D12 made the duplicate harmless, not absent). And, as for equal cells,
nothing under \code{src/} was touched: \code{farBlockX!} replaces the \code{egoFunOut!}
loop of \code{genEgoCrcExt!} only once someone wires it there.

\paragraph{The contact path is the remaining bottleneck, and it is now also the remaining
accuracy problem.} After this work a $128^3$ build at $\lambda/32$ is $6.73$~s on twelve
threads of which $5.14$~s is \code{wekTrp} at \code{intOrd 48}
(Section~\ref{sec:bch}); and Section~\ref{sec:ver:pos} shows that the same integrals are
wrong by $117\,\%$ on a slender touching shell. The moment series of the companion document
computes them exactly. \emph{Both problems have the same fix, and it is not in this
deliverable.}
